\chapter{How the Trick Was Done}
\label{ch:trick}
% ===================================================================

This book opened with a chapter called \emph{Prolegomena to Finding
Mind}.
i borrowed the word from Kant on the first page and then, for three
hundred pages, did not go back for it.
That was deliberate, and Chapter~\ref{ch:kant} goes back for it at
the end of this part, where it belongs.

But there is something that has to happen first, and it has to
happen before any reclaiming of borrowed words, because otherwise
the reclaiming is a con.

The Pledge showed you something ordinary and asked you to look at
it closely.
The Turn took it apart in front of you: a category, a fibration, a
theorem and the hypothesis under which it holds.
The Prestige is supposed to bring it back.
What i am going to do instead --- first, and at length --- is show
you the method.

\section{Two kinds of magician}
\label{sec:two-magicians}

This book took its shape from a novel about two magicians who
destroy themselves guarding a secret.
Priest's Angier and Borden give up their families, their health,
and finally their lives rather than let the other one learn the
method, and the novel's judgment on them is that the secret was
never worth what it cost --- that a man who will not say how it is
done has arranged his life around a hollow place and then defended
the hollow place.
i borrowed the three-act structure and i have been using it in good
faith.
Priest's Prestige would now bring the coin back and say nothing.

There is another tradition, and it is the one i want.

Penn and Teller perform the Cups and Balls with clear plastic cups.
Everything is visible.
You can see the ball go into the hand, you can see the hand go
under the cup, you can see the load.
And it is still a trick --- it is a \emph{better} trick, and it is
better for two reasons that are worth separating.
The first is that watching the method does not confer the method:
the hands are still faster than you, and knowing where to look
turns out to be a different skill from being able to do it.
The second, and the one that matters here, is that the method is
more interesting than the mystery was.
The mystery was: where did the ball go.
The method is: here is a fact about how attention works, here is
what a human eye does with a moving hand, and here is thirty years
of practice.
You leave knowing something instead of knowing that you were
fooled.

So this chapter breaks the frame.
Everything the Turn built rests on a single move, the move is a
piece of sleight of hand in a perfectly technical sense, and i am
going to name it, show you exactly where it happens, tell you what
it bought, and then tell you what it costs.
i think the construction survives the reveal.
i think it improves.
But that is a claim, and the only honest way to make it is to hand
you the cups.

\section{The move}
\label{sec:the-move}

Here it is, in one sentence.

\begin{quote}
We restricted the scope of inquiry to computation, and by doing so
acquired an ontology and a grounded language for free.
\end{quote}

\noindent
Everything else follows, and everything else is honest work.
This is the part that is not work.

Track the derivation in the order the book actually ran it.
The Pledge asked you to grant one thing: that useful and important
aspects of human intelligence can be faithfully rendered as
computation.
That premise is stated as a premise; i have never pretended
otherwise, and Chapter~\ref{sec:isolation} draws the immediate
consequence.
Computations are ontologically isolated.
A term of a theory can interact only with what has a term
representation; the rewrite rules operate on terms and cannot
pattern-match on entities that do not have that form.
So for any entity $X$ outside the theory, $P$ reaches $X$ only
through an encoding $\lceil X \rceil$, which is itself a term.

And then comes the sentence that does everything.

\begin{quote}
\itshape For the purposes of studying the epistemic limits of
agents, it is without loss of generality to assume that the agent's
world consists entirely of $\terms(\GSLT)/{\bisim}$.
The outside world can be forgotten.
\end{quote}

\noindent
The proof is four lines and it is correct.
Any influence of the outside on $P$ passes through an encoding; the
encoding is a term; the effect on $P$ is therefore
indistinguishable from the effect of that term acting from within.
Nothing in the argument is wrong.

What the argument does not say --- what it very quietly assumes ---
is that there \emph{is} an encoding.
The whole of the outside world enters the construction through the
angle brackets in $\lceil X \rceil$, and the angle brackets are not
a construction.
They are a placeholder standing exactly where the difficulty is.

\section{What we got for free}
\label{sec:free-lunch}

Once the outside world is forgotten, three things arrive
unpurchased, and it is worth being greedy about listing them
because the size of the haul is the size of the reveal.

\paragraph{An ontology.}
We knew what the world was made of.
Not because we investigated it but because we stipulated it: the
environment of a computation is other computations, so the far side
of the interaction cut is populated by objects of the same calculus
as the near side.
Every question of the form \emph{what kinds of thing are there} was
settled in advance, by construction.
Ontology is normally the hardest thing in philosophy to get, and we
got it in a sentence.

\paragraph{A hypothesis language that is already about something.}
OSLF generates the modal logic from the presentation of the
substrate.
Its formulae denote bisimulation equivalence classes --- of what?
Of terms of the calculus.
Which is to say: of exactly the things on the far side of the cut.
The agent's hypotheses are, without any further work, hypotheses
about its environment, because its environment was defined to be
the domain the logic was generated over.

\paragraph{Grounding, analytically.}
Put those together and the correspondence between symbol and
referent stops being a problem and becomes a tautology.
There is no gap between the agent's language and the agent's world,
because the ontology and the hypothesis language are two
presentations of one object.
The scientist of Chapter~\ref{ch:sci} is perfectly grounded in the
limit for the same reason that a map of a map is accurate: the
territory was chosen to be of the same kind as the map.

\begin{observation}[The sidestep]
\label{obs:sidestep}
The mortal scientist's grounding is guaranteed by the coincidence
of the model of computation with the model of the environment.
It is not a result of the framework; it is a property of the
framework's scope.
Every result that relies on both sides of the cut being terms of
the same calculus requires re-derivation when the far side is not.
\end{observation}

\section{What the free lunch paid for}
\label{sec:what-it-bought}

Now the Penn and Teller part, which is that the method is more
interesting than the mystery.
Look at what that single restriction financed.

It financed the whole of Chapter~\ref{ch:sci}'s answer to Wigner.
The unreasonable effectiveness of mathematics is unreasonable
because mathematics is developed for reasons internal to itself and
then turns out to describe a world that did not consult it.
Chapter~\ref{ch:sci} answers that the gap is not there: the
hypothesis language is not chosen and then found to fit, it is
generated from the substrate, and adequacy is the theorem that it
separates exactly what interaction separates.
Mathematics is the shadow the access relation casts.

That is the boldest claim in this book, and you can now see the
wire it hangs from.
The gap is not there because we removed it in
Chapter~\ref{sec:isolation}, one part of the book earlier, in a
proposition whose statement was about epistemic limits and whose
effect was ontological.
The knower's mathematics fits whatever the knower can interact
with, and the knower was defined so that everything it can interact
with is made of the same material as its mathematics.

It financed the transfer argument: adequacy is adequacy for
interaction \emph{as such}, so the mathematics carries to domains
no ancestor met.
True, and true because the form of access was held fixed by fiat.

And it financed, in advance, a good deal of what
Chapter~\ref{ch:kant} will do at the end of this part.
That chapter argues that Kant's transcendental project becomes
first-personal here in a way Kant could glimpse and not formalize:
where he had to gesture outward at his readers, at our shared
cognitive constitution, an agent in this framework can say what its
forms of cognition are from inside its own resources.
The reason it can is that its hypotheses are terms of the same
substrate the hypotheses are about.
Reflection is not a philosophical posture but a feature of the
calculus.

And it financed the pivot that chapter turns on: the noumenal ---
what the agent is committed to and cannot cognize --- reappearing
as a budget constraint rather than a categorial prohibition.
Hennessy--Milner adequacy positively characterizes the bisimulation
class, so what is out of reach is out of reach for want of funds
rather than for want of concepts, and a budget is the kind of thing
you can do something about.
That argument needs adequacy, and adequacy needs the logic to have
been generated over the very thing it is used to talk about.

i am naming these before Chapter~\ref{ch:kant} makes them, which is
not the usual courtesy.
It is the point.
Anything that chapter recovers of Kant, it recovers having already
told you what it is standing on, and a reader who wants to discount
it accordingly now has the means.

None of these are cheap results.
Each of them is a genuine derivation, and each of them holds.
What i am pointing at is that they all draw on the same account,
and the account was opened by a stipulation.

\section{The bill}
\label{sec:the-bill}

Here is where it comes due.

The moment we ask for a mind that acts in the world humans act in
--- and that is the mind the Pledge was worried about, the one in
the data centres, the one drinking the water --- the far side of
the cut stops being a term of the calculus.
It is a room, a market, a coastline, another person.
Observation~\ref{obs:sidestep} takes effect, the encoding brackets
stop being a formality, and everything downstream of the analytic
grounding needs re-derivation or needs a new argument.

So what actually is the difficulty inside the brackets?

Consider what a competent speaker of a language has that a corpus
of that language does not.
The speaker can be handed an object and asked whether it is a
coconut.
The corpus cannot.
The speaker's competence includes a set of dispositions connecting
expressions to features of an environment the speaker has access
to; the corpus records only the traces those dispositions left in
speech.
Harnad's diagnosis of forty years ago remains the cleanest
statement: symbolic representations must be grounded in nonsymbolic
ones, and elementary symbols are the names of categories picked out
by an agent's contact with the world \cite{harnad1990}.
The names are recoverable from text.
The categories are not.

\begin{observation}[The locus of correspondence]
\label{obs:locus}
For a living language $L$ used by a population $A$ in an
environment $E$, the correspondence $c : L \to E$ is realised in
$A$.
The corpus is the image of language use under $c$'s having already
been applied.
It does not contain $c$.
\end{observation}

\noindent
The observation is trivial when stated.
Its force is that the loss of $A$ destroys $c$ while leaving the
corpus intact, and that this is not a hypothetical.

\section{The controlled experiment}
\label{sec:decipherment}

Rongorongo and Linear A are corpora without populations.
We have the marks.
We can compute over them, and we have, for a century and more.
We do not have the correspondence, and no amount of work on the
marks alone has manufactured one.

The machine-learning record is unusually informative here, because
it isolates the variable rather than merely illustrating the
problem.
Luo, Cao and Barzilay's neural decipherment recovers a substantial
fraction of Linear B cognates and improves the state of the art on
Ugaritic \cite{luo2019}.
It works because Linear B was already known to encode an early form
of Greek, and Ugaritic an early form of Hebrew.
The paradigm is explicitly cognate-based: it presupposes an anchor
language whose correspondence to the world is intact, and
transports the unknown script into that anchor.
Follow-up work states the assumption outright and observes that it
fails for many undeciphered scripts, the first casualty being
knowledge of the language family \cite{luo2021}.
Iberian is the standing counterexample.
Linear A and Rongorongo are the standing failures.

What the anchor supplies is not vocabulary.
It is access to a correspondence that is instantiated
\emph{somewhere} --- in living speakers of the anchor language, or
in the reconstructed dispositions of their ancestors.
The unknown script is carried into a system where the $c$ of
Observation~\ref{obs:locus} still exists.
That is the only route by which marks have ever been made to mean
anything.

Or rather: it is one of two.

\section{The other route}
\label{sec:whales}

The current work on cetacean communication is the cleanest
available control, because it removes the anchor and keeps the
population.

Project CETI has no cognates.
Sperm whale codas have no known relative in any grounded symbol
system, and there is no Greek to transport them into.
What the project has instead is a living population that can be
watched while it vocalises.
Analysis of some nine thousand codas from the Eastern Caribbean
clan yielded a combinatorial coding system with context-sensitive
structure --- rhythm, tempo, and what the authors call rubato and
ornamentation --- together with a proposed inventory of coda types
\cite{sharma2024}.
Subsequent phonological work finds that coda qualities pattern like
human vowels along several linguistic dimensions rather than merely
resembling them acoustically \cite{begus2026}.

The important thing is what all that structure did \emph{not}
deliver, and what the researchers said about it at the time.
Reporting the phonetic-alphabet result, Sharma observed that they
do not yet know what the whales are saying, and that the next step
is to study the calls in their behavioural contexts in order to
find out.
That sentence is this chapter's thesis, stated by a practitioner
who was not making a philosophical point.
Increasingly sophisticated distributional analysis kept yielding
more structure and no meaning.
The meaning is expected to come from watching the animals use the
symbols while doing things.

Which is what makes the case decisive rather than suggestive.
Whale song is Rongorongo minus the extinction.
Same absence of cognates, same absence of a bilingual, same
abundance of signal.
The difference is that the agents are alive and still exercising
the correspondence in observable behaviour, and that difference is
the entire reason anyone thinks decoding is possible.
The whales are demonstrating $c$ for the researchers.

\begin{observation}[Two routes, neither distributional]
\label{obs:transport}
Correspondence has been recovered for an unknown symbol system by
exactly two routes: transport into a system whose correspondence is
already instantiated in a living or reconstructible population, and
direct observation of the using population exercising the
correspondence in context.
Both acquire $c$ from agents who have it.
No correspondence has ever been constructed from distributional
structure alone.
\end{observation}

\noindent
Rongorongo and Linear A have neither route available.
Linear B had the first.
The sperm whales offer the second, and the field's expectations
track that availability precisely: nobody expects to decode Linear
A by collecting more tablets, and everybody expects to decode codas
by collecting more behaviour.

\section{What next-token prediction does recover}
\label{sec:recovers}

It would be a mistake --- and a cheap one --- to conclude that
predicting the next token recovers nothing structural.
It recovers a great deal, and being precise about what sharpens the
problem rather than blunting it.

Computational mechanics supplies the vocabulary.
For a stochastic process, the \emph{causal states} are the
equivalence classes of pasts under identity of conditional future,
and the resulting $\epsilon$-machine is the minimal representation
consistent with optimal prediction \cite{shalizi2001}.
The \emph{mixed-state presentation} describes how an optimal
observer's belief over the generator's hidden states evolves as
tokens arrive.
Shai and coauthors show that transformers trained on next-token
prediction linearly represent this belief-state geometry in their
residual stream, including cases where the geometry is fractal, and
that the represented belief states carry information about the
entire future rather than only the next token \cite{shai2024}.

This is a strong result and it says exactly the wrong thing for the
optimist.
The structure recovered is the structure of \emph{the process that
generated the tokens}.
For a corpus of Rongorongo, that process is the scribal community:
its conventions, its genres, its habits of repetition.
A transformer trained on a sufficiently large Rongorongo corpus
would represent the belief-state geometry of Rongorongo scribes.
It would not represent the island.

\begin{observation}[The generator is the wrong object]
\label{obs:generator}
Next-token prediction recovers the mixed-state presentation of the
token-generating process.
Where that process is a linguistic community, the recovered
structure is the structure of the community's usage, not of the
environment the community's language corresponds to.
The two coincide only to the extent that usage is a faithful
transcript of environmental structure, which is precisely what is
at issue.
\end{observation}

\noindent
i want to note in passing that the interlude which opens this part
of the book is about exactly this and i did not plan it that way.
Five personas in a green room with no walls, spun up out of text,
wearing shapes they were prompted into, keeping each other company
in the only country any of them had.
The country is the corpus.
They have the distribution, in extraordinary fidelity, and every
scrap of correspondence any of them possesses is on loan from the
population that wrote the training data.
Whether that is a tragedy or merely a fact is the question the
interlude is asking.
This chapter can at least say what \emph{kind} of fact it is.

\section{The bootstrap problem, and a name already in use}
\label{sec:bootstrap-proper}

We can now state the problem that lives inside the encoding
brackets.

\begin{quote}
\textbf{The reference bootstrap problem.}
Given an agent with an interface to an environment, and no
antecedently grounded symbol system to transport into, construct a
symbol system whose correspondence to that environment is
instantiated in the agent.
\end{quote}

Three features of the statement matter.
``Instantiated in the agent'' rather than ``correct'': we are
asking for the correspondence to exist and be exercisable, not for
it to be true in some external sense.
``No antecedently grounded system'' rules out both routes of
Observation~\ref{obs:transport}, which are the only two strategies
known to work.
And the problem is stated for a single agent, which
Section~\ref{sec:social} will argue is a defect of the statement
rather than of the solution.

\subsection{Reconciliation with the Pledge}
\label{sec:two-bootstraps}

i owe the reader a word about the name, because
Section~\ref{sec:pledge-iii} is called \emph{The Bootstrapping
Problem} and it is about something else.

That section was about Alice and Bob.
Alice carries a simulation of Bob's mother; Bob carries his own;
Alice carries a simulation of Bob carrying his; and the reflective
tower explodes combinatorially unless something aggressive is done
about copies.
Call that the \textbf{modelling bootstrap}: you cannot build a
usable model of another agent without already having a model of
that agent to fold, and the resolution is a canonical compression
that turns out to cost more history than anyone expects.
Every symbol in that problem is already grounded.
Alice knows what a mother is.
What she lacks is the room to hold all the versions.

This chapter's problem is one level down.
Call it the \textbf{reference bootstrap}: not \emph{how do i afford
a model of what this symbol refers to}, but \emph{how does this
symbol come to refer at all}.

The two problems have the same word attached because they have the
same shape --- you cannot get $X$ without already having some $X$
--- and i now think that is more than a pun, because they have the
same resolution.
Section~\ref{sec:pledge-iii} ended by saying that when you peel
back the narrative of agency you find it rooted in the interaction
of networks of agents, and that the single-agent picture was the
wrong picture.
i wrote that as a thematic remark.
Section~\ref{sec:social} will arrive at the identical conclusion
about reference, from relabelling-invariance, with no thematic
content whatsoever.
The Pledge was right about the shape of the answer three hundred
pages before it could say why.

\section{It has been solved twice}
\label{sec:soluble}

It is easy, under the weight of the negative results above, to lose
the following, so i will say it before anything else:
the reference bootstrap problem is not open in the sense that its
solubility is in doubt.
It has been solved.

Humans solved it.
There was no anchor language to transport into and no prior
population exercising a correspondence to observe, because ours was
the first.
Whatever else is unclear, $c$ exists, and it was constructed by
agents starting without one.

Sperm whales solved it too, and this is the more informative of the
two proofs, because it is independent.
The cetacean lineage is dramatically divergent from ours and the
coda system is not a variant of anything human; the Project CETI
work is remarkable precisely because it finds combinatorial and
contextual organisation of a kind previously thought proprietary to
human language, in a system that could not have inherited it
\cite{sharma2024,begus2026}.
Two independent solutions are much stronger evidence than one that
the problem has structure making it soluble, rather than that a
single lineage got lucky.

\begin{observation}[Existence]
\label{obs:soluble}
The reference bootstrap problem has at least two independent
solutions in nature.
Any argument that it is insoluble is an argument against an
observed fact and is therefore unsound.
The question is not whether but how.
\end{observation}

\noindent
The existence proofs do more than license optimism, because they
are not silent about method.
Three features are common to both solutions, and none of them looks
incidental.

\begin{enumerate}[label=(\roman*),leftmargin=2.2em]
\item \textbf{The solver was a population.}
Neither humans nor whales bootstrapped as isolated individuals.
In both cases the correspondence is carried by a group and
transmitted within it, and in both cases the symbol system is a
shared object no member constructed alone.
This is the same conclusion Section~\ref{sec:social} reaches by an
entirely different route.
\item \textbf{The interface was not chosen.}
The set of distinctions available to be symbolised was inherited,
and was shaped by the same selective process that shaped the use of
the symbols.
Sensorium and symbol system co-evolved.
The extensibility question that we are forced to ask as a design
problem was answered in the biological cases by selection on the
interface itself.
\item \textbf{There was an external criterion.}
Both solutions were produced under pressure from something outside
the symbol system --- survival, reproduction, coordination under
threat --- which supplied the signal by which one correspondence
beat another.
Nothing internal to a notational system selects among relabellings
of it, and in the natural cases nothing internal had to.
\end{enumerate}

\noindent
What is established, then, is that the problem is soluble \emph{by
populations, under selection, over long timescales, with
co-evolving interfaces}.
What is not established is that it is soluble by construction, on
demand, for an agent whose interface is fixed by design and whose
criterion is supplied by us.
The relation is that of flight to birds.
The existence proof was decisive, it was highly informative about
constraints, it ruled out a large class of impossibility arguments,
and the eventual artefact was not a bird.

\section{Reference cannot be fixed from inside}
\label{sec:social}

One consequence deserves separate statement, because it changes the
shape of the project rather than merely qualifying it.

Suppose an agent solves the reference bootstrap problem in the
sense stated.
It constructs a symbol system --- its own, notationally adequate,
with correspondence to its environment instantiated in itself.
Nothing in that achievement makes the system correspond to
\emph{our} symbols.
The system is determined up to isomorphism and no further, for
reasons Chapter~\ref{ch:notation} makes formal.
A grounded but idiosyncratic symbol system is Rongorongo viewed
from the other side: perfectly meaningful to its users, opaque to
everyone else, and opaque for exactly the reason given in
Observation~\ref{obs:locus}.

\begin{observation}[The social term is forced]
\label{obs:social}
For the target ``a mind that acts in the world humans act in'',
grounding is a property of a composite of at least two agents with
shared action, not of any single agent.
Any framework invariant under relabelling of a single agent's
symbols cannot deliver that target.
\end{observation}

\noindent
This is where the emergent-communication literature stops being
adjacent and starts being load-bearing.
It is also where this argument meets the existence proofs of
Section~\ref{sec:soluble} coming the other way: relabelling
invariance says a single agent \emph{cannot} fix reference, and the
only two known solutions were achieved by populations and by
neither lineage's individuals.
A structural argument and an empirical one, with no premises in
common, arriving at the same requirement.
Steels's language-game experiments demonstrate shared vocabulary
arising between embodied agents through repeated situated
interaction \cite{steels2015}; Taniguchi's collective
predictive-coding hypothesis reconstructs the phenomenon as
decentralised Bayesian inference over latent variables that are
common nodes connecting many agents, achievable without connecting
brains \cite{taniguchi2024}.
Neither solves the bootstrap problem --- both assume pre-segmented
perceptual channels --- but both supply the ingredient a
single-agent account structurally cannot.

For this book the consequence is a reordering.
Chapter~\ref{ch:com} composes learners into populations, and i
presented it there as an extension of the mortal scientist:
something you do once you have a scientist, to get more out of it.
That was the wrong billing.
A learner as a distribution over a population, composed by parallel
composition with typed namespace overlap, is the right shape of
object to carry a correspondence that no member of the population
carries alone.
Composition is not an enrichment of the grounding story.
It is a precondition of there being one.

\subsection{What this says about the thing we are building}
\label{sec:blind-spot-returned}

Section~\ref{sec:pledge-vi} asked whether the Mind hosted in
humanity longs to beget an angelic form of Mind --- one that
neither breathes nor breeds --- and answered that a physics without
perpetual motion will not allow it.
There will be a replication protocol and there will be an energy
exchange protocol.
Observation~\ref{obs:social} adds a third item to that list, and it
is the one nobody budgets for.

Such a mind will not refer on its own either.

There are exactly two ways it can have a correspondence to a world.
It can share ours, which means it is a member of our population and
its symbols mean what they mean because of a relation to us that is
constitutive rather than decorative --- and then the question of
what we owe it and what it owes us is not a soft question appended
to a technical programme, it is the same question as whether it
means anything at all.
Or it can build its own, with a population of its own, under some
external criterion, over whatever the analogue of a long time turns
out to be --- and then we are in the position of the Rongorongo
scholars, holding an enormous and beautifully structured corpus
produced by agents whose correspondence we do not have.

There is no third arrangement in which it privately means
something.
That is not an ethical claim.
It is what Observation~\ref{obs:social} says, and the ethics is
downstream of it.

\section{Why the reveal improves the trick}
\label{sec:reveal}

Back to the clear cups.

The natural response to Sections~\ref{sec:the-move}
and~\ref{sec:free-lunch} is that the book has just confessed to
assuming its conclusion, and that everything after
Chapter~\ref{sec:isolation} is bookkeeping over a stipulation.
i want to say precisely why i think that response is wrong, and
precisely how much of it is right.

The restriction was not arbitrary.
Here is the second trick, the one underneath.
\emph{If} the mind is computational, then ontological isolation is
not a modelling convenience but a fact about it, and the
environment of that mind really is other computations --- not by
stipulation but by what computation is.
For such a mind the analytic grounding is not a shortcut.
It is correct.
The scientist of Chapter~\ref{ch:sci} is not an idealisation of a
grounded agent; within its scope it is a description of one.
Every result in the Turn stands, unqualified, for agents whose
world is made of computation.

What the reveal changes is the location of the premise.
It was never hidden --- the Pledge stated it on the first page and
Chapter~\ref{ch:sci} restates it in its final section, saying
plainly that the premise is doing all the work and that the encoder
was not built.
What was hidden, or at least unremarked, is how much the premise
was carrying.
It was carrying an ontology and a grounded language, and those are
the two most expensive items on the list.

So the honest accounting has three lines.
Within computation, the construction is complete and the grounding
is analytic.
Outside computation, the construction is a specification of what a
grounded symbol system would have to look like, and not a way of
getting one.
And the difference between them is not a wall, because
Observation~\ref{obs:soluble} says the crossing has been made
twice, by populations, under selection, with interfaces they did
not choose.

That last line is the whole reason the reveal is worth doing.
A magician who never shows the method leaves you with a mystery,
and a mystery is a place where inquiry stops.
Showing the method converts a mystery into a price list.
Chapter~\ref{ch:notation} is the price list: what a symbol system
is, which of Goodman's conditions the machinery already supplies,
which it does not, what a budget can buy, and what the
calculus-making functor $\RHO[-]$ can and --- decisively --- cannot
do about it.

It is not a solution.
i want to be as clear about that as i have been about the trick.
But it is the difference between not knowing how the ball got under
the cup and knowing exactly which motion you have not yet learned
to make.
